\DocumentMetadata{
  pdfversion=2.0,
  pdfstandard=ua-2,
  lang=en-US,
  tagging=on,
  tagging-setup={math/setup=mathml-SE,tabsorder=structure}
}
\documentclass[11pt,letterpaper]{article}
\usepackage[margin=0.68in,includefoot]{geometry}
\usepackage{newcomputermodern}
\usepackage{amsmath,amscd,mathtools}
\usepackage{tikz,adjustbox}
\providecommand{\square}{\mdlgwhtsquare}
\usepackage[hidelinks]{hyperref}
\hypersetup{
  pdftitle={Analysis PhD Exam, May 2016},
  pdfauthor={University of Florida Department of Mathematics},
  pdfsubject={Graduate mathematics examination},
  pdfdisplaydoctitle=true,
  bookmarksopen=true
}
\setcounter{secnumdepth}{0}
\setlength{\parindent}{0pt}
\setlength{\parskip}{0.28em}
\setlength{\emergencystretch}{3em}
\setlength{\leftmargini}{2.15em}
\setlength{\leftmarginii}{2.65em}
\setlength{\leftmarginiii}{3em}
\renewcommand{\labelenumi}{\textbf{\theenumi.}}
\pagestyle{empty}
\raggedbottom
\allowdisplaybreaks
\newenvironment{examproblems}[1]{%
  \begin{enumerate}%
  \setcounter{enumi}{#1}%
  \setlength{\topsep}{0.35em}%
  \setlength{\partopsep}{0pt}%
  \setlength{\itemsep}{0.48em}%
  \setlength{\parsep}{0.18em}%
}{\end{enumerate}}
\newenvironment{examsubparts}{%
  \begin{enumerate}%
  \setlength{\topsep}{0.18em}%
  \setlength{\partopsep}{0pt}%
  \setlength{\itemsep}{0.16em}%
  \setlength{\parsep}{0.1em}%
}{\end{enumerate}}
\newenvironment{examinstructions}{%
  \par\begingroup\itshape
}{%
  \par\endgroup\vspace{0.18em}
}
\makeatletter
\renewcommand\section{\@startsection{section}{1}{0pt}%
  {-1.25ex plus -.25ex minus -.1ex}%
  {0.55ex plus .1ex}%
  {\normalfont\large\bfseries}}
\renewcommand{\@maketitle}{%
  \newpage\null\vskip -1em%
  {\footnotesize\bfseries
    \href{https://www.ufl.edu/}{University of Florida}, \href{https://math.ufl.edu/}{Department of Mathematics}\par}%
  \vskip -0.5em%
  \tagstructbegin{tag=Title}%
    {\LARGE\bfseries\href{https://gma.math.ufl.edu/exams/analysis-phd/}{Analysis PhD Exam}, May 2016\par}%
  \tagstructend%
  \vskip 0.25em%
  \tagmcbegin{artifact}%
    \hrule height 1pt%
  \tagmcend%
  \vskip 1em%
}
\makeatother
\title{Analysis PhD Exam, May 2016}
\author{}
\date{}
\begin{document}
\maketitle
\thispagestyle{empty}
\begin{examinstructions}
Do THREE problems from Part A and THREE problems from Part B. Write solutions in a neat and logical fashion, giving complete reasons for all steps.
\end{examinstructions}
\section{Part A}
\begin{examproblems}{0}
\item
Consider Lebesgue measure on the real line~\(\mathbb{R}\). Give an example for each of the following, if possible. If not possible, give a brief explanation. A sequence \((f_n)\) in \(L^1(\mathbb{R})\) converging to an~\(f\) in \(L^1(\mathbb{R})\)...
\begin{examsubparts}
\item[\textnormal{a)}]
...in the \(L^1\) norm but not in measure,
\item[\textnormal{b)}]
...in measure but not in the \(L^1\) norm,
\item[\textnormal{c)}]
...in the \(L^1\) norm but not a.e.,
\item[\textnormal{d)}]
...a.e. but not in measure.
\end{examsubparts}
\item
Let~\(\nu\) be a signed measure on \((X,\mathcal{M})\).
\begin{examsubparts}
\item[\textnormal{a)}]
Define what it means for a set to be positive, negative or null for~\(\nu\), and state the Hahn decomposition theorem.
\item[\textnormal{b)}]
Prove there exist unique positive measures \(\nu^+,\nu^-\) on~\(\mathcal{M}\) so that \(\nu=\nu^+-\nu^-\) and \(\nu^+\perp\nu^-\).
\end{examsubparts}
\item
Let~\(\mu\) be a finite, regular Borel measure on \([0,1]\). Suppose that 
\[
\int_0^1 x^n\,d\mu=0 \quad \text{for } n=0,1,2,\ldots
\]
 Prove that \(\mu=0\).
\item
Let \((X,\mathcal{M},\mu)\) be a~\(\sigma\)-finite measure space. Let~\(\mathcal{N}\) be a sub-\(\sigma\) algebra of~\(\mathcal{M}\) and let~\(\nu\) be the restriction of~\(\mu\) to~\(\mathcal{N}\). Prove that for each \(f\in L^1(\mu)\), there exists a \(g\in L^1(\nu)\) such that 
\[
\int_E f\,d\mu=\int_E g\,d\nu
\]
 for all \(E\in\mathcal{N}\).
\end{examproblems}
\section{Part B}
\begin{examproblems}{0}
\item
Let~\(\mathcal{X}\) be a normed vector space. Say a sequence \((x_n)\) from~\(\mathcal{X}\) converges weakly to \(x\in\mathcal{X}\) if \(f(x_n)\to f(x)\) for all \(f\in\mathcal{X}^*\). Prove that if~\(\mathcal{M}\) is a norm closed subspace of~\(\mathcal{X}\), and \((x_n)\) is a sequence in~\(\mathcal{M}\) converging weakly to \(x\in\mathcal{X}\), then \(x\in\mathcal{M}\).
\item
Let~\(\mathcal{H}\) be a Hilbert space and \(\mathcal{M},\mathcal{N}\) closed subspaces of~\(\mathcal{H}\). Prove that if \(\mathcal{M}\perp\mathcal{N}\), then \(\mathcal{M}+\mathcal{N}\) is closed.
\item
Fix \(1<p<\infty\), let~\(\mu\) be a~\(\sigma\)-finite measure, and let \((f_n)\) be a sequence in \(L^p(\mu)\). Suppose that there exists a function \(f\in L^p(\mu)\) such that 
\[
\int f_ng\,d\mu\to\int fg\,d\mu
\]
 for every \(g\in L^q(\mu)\), where \(\frac1p+\frac1q=1\).
\begin{examsubparts}
\item[\textnormal{a)}]
Prove that \(\sup_n\{\lVert f_n\rVert_p\}<+\infty\).
\item[\textnormal{b)}]
Prove that \(\lVert f\rVert_p\leq\limsup\lVert f_n\rVert_p\).
\end{examsubparts}
\item
Let \((\varphi_n)\) be a sequence in \(L^1(\mathbb{R})\) with the following properties:

\par
Prove that for every \(f\in L^1(\mathbb{R})\), \(f*\varphi_n\to f\) in \(L^1(\mathbb{R})\).
\begin{examsubparts}
\item[\textnormal{i)}]
\(\varphi_n\geq 0\) for all~\(n\),
\item[\textnormal{ii)}]
\(\displaystyle \int_{\mathbb{R}}\varphi_n(x)\,dx=1\) for all~\(n\), and
\item[\textnormal{iii)}]
for every \(\delta>0\), \(\displaystyle \lim_{n\to\infty}\int_{|x|>\delta}\varphi_n(x)\,dx=0\).
\end{examsubparts}
\end{examproblems}
\end{document}
